\section{Future Work}

To build upon the findings and address the identified limitations, several research directions are planned:
\begin{itemize}
    \item \textbf{Experimental Hardware Validation:} Testing and validating the PINN digital twin on physical lithium-ion battery cells operating inside a thermal climate chamber under dynamic load profiles.
    \item \textbf{Multi-Node Physics Formulation:} Extending the physics loss to a multi-state (core and surface) thermal formulation to resolve core-surface phase lags without relying solely on softplus parameter bounds.
    \item \textbf{Comprehensive Baseline Benchmarking:} Expanding model comparisons to include Recurrent Neural Networks (LSTM/GRU), Temporal Convolutional Networks, Transformers, and continuous-time Neural ODEs.
    \item \textbf{Multi-Seed Statistical Analysis:} Conducting multi-seed training runs to compute mean, standard deviation, and statistical confidence intervals for rollout performance metrics.
    \item \textbf{Online Parameter Adaptation:} Integrating recursive estimation frameworks (e.g., Extended Kalman Filtering) for real-time online updates of $hA$ and $m C_p$ as cooling conditions and flow rates vary.
    \item \textbf{Aging and SOH Integration:} Incorporating capacity fade and internal resistance growth into the loss function to model thermal behavior changes across the battery lifecycle.
    \item \textbf{Embedded Edge Deployment:} Exporting trained models via ONNX Runtime or C++ TensorRT to benchmark execution speed and memory utilization on automotive microcontrollers.
\end{itemize}
